\section{Computational Cost}
\label{appendix:cost}

\begin{table}[ht]
\centering
\begin{tabular}{lccm{5cm}}
\toprule
\textbf{Method} & \textbf{\#LLM Calls} & \textbf{Latency (Mean)} & \textbf{System Description} \\
\midrule
Single Agent & 1 & 3.3 mins & Single monolithic prompt and generation. \\
\addlinespace
AI Scientist-v2 & $\sim$40--45 & 35.1 mins & Multi-step pipeline (20-round citation scan, drafting, $3\times$ reflections, and vision critique). \\
\addlinespace
\ourmethod{} & $\sim$60--70 & 39.6 mins & Multi-agent orchestration (outline generation, parallel literature search, sequential citation verification, writing, plotting with critique, and a $3\times$ content refinement loop). \\
\bottomrule
\end{tabular}
\caption{Comparison of computational cost and single-paper latency. This table outlines the estimated number of LLM API calls, mean execution time, and core architectural components for each pipeline.}
\label{tab:runtime_efficiency}
\end{table}

Table~\ref{tab:runtime_efficiency} compares the computational cost and single-paper latency of \ourmethod{} against the Single Agent and AI Scientist-v2 baselines. Although the Single Agent executes rapidly, it fails to produce the rigorous citations and data-grounded visuals generated by our system. Despite requiring more LLM calls ($\sim$60--70) than AI Scientist-v2 ($\sim$40--45), \ourmethod{} maintains a highly competitive mean processing time of 39.6 minutes (compared to 35.1 minutes for AI Scientist-v2). To calculate these mean latency values, we evaluated all methods on a random subset of 10 papers (5 from CVPR and 5 from ICLR) using a single worker. For each method, we removed the highest and lowest execution times to mitigate the impact of API network anomalies before computing the final average.

The efficiency of \ourmethod{} is achieved by strategically decoupling the paper discovery and verification pipeline. To optimize throughput, we execute retrieval in two distinct stages: (1) \textbf{Parallel Candidate Discovery}, which leverages 10 concurrent workers for search-grounded LLM calls to rapidly pool candidate papers; and (2) \textbf{Sequential Citation Verification}, which safely processes the pooled candidates through Semantic Scholar at the maximum allowable rate (1 query per second). This architecture successfully combines the high-concurrency tolerance of the LLM API with the strict throughput limits of the Semantic Scholar API to prevent quota-induced latency.

The LLM calls are distributed across our multi-agent system as follows:
\begin{itemize}
    \item \textbf{Outline Agent} (1 call): Initial manuscript structuring and planning.
    \item \textbf{Hybrid Literature Agent} ($\sim$20--30 calls): Execution of the decoupled paper discovery and citation verification pipeline described above.
    \item \textbf{Plotting Agent} ($\sim$20--30 calls): Few-shot retrieval, visual planning, image generation, VLM-guided critique and redraw cycles, and context-aware captioning.
    \item \textbf{Section Writing Agent} (1 call): A single, comprehensive multimodal call to draft and compile the complete \mylatex{} manuscript.
    \item \textbf{Content Refinement Agent} ($\sim$5--7 calls): Score-driven, iterative reflection and manuscript revision.
\end{itemize}

Ultimately, this computational investment is well justified. \ourmethod{} delivers substantial improvements in overall manuscript quality and citation depth, while adding only minimal latency compared to the AI Scientist-v2 baseline.